\documentclass{aastex631}

\usepackage{amsmath,amssymb}
\usepackage{natbib}

\newcommand{\dd}{\mathrm{d}}
\newcommand{\pp}{\partial}
\newcommand{\Hp}{H_P}
\newcommand{\Cp}{C_P}
\newcommand{\Tint}{T_{\rm int}}
\newcommand{\Teff}{T_{\rm eff}}
\newcommand{\Mjup}{M_{\rm Jup}}
\newcommand{\Mearth}{M_\oplus}
\newcommand{\softplus}{\mathrm{softplus}}
\newcommand{\Vmlt}{v_{\rm MLT}}
\newcommand{\VNR}{v_{\rm NR}}
\newcommand{\VR}{v_{\rm R}}
\newcommand{\Vconv}{U_{\rm conv}}
\newcommand{\Roo}{\mathrm{Ro}}
\newcommand{\Rofunc}{R(\Roo_{\rm NR})}
\newcommand{\NCsq}{N_C^2}
\newcommand{\Dmlt}{D_{\rm MLT}}
\newcommand{\Dent}{D_{\rm ent}}
\newcommand{\Dov}{D_{\rm ov}}
\newcommand{\Dmicro}{D_{\rm micro}}
\newcommand{\sigmaconv}{\sigma_{\rm conv}}
\newcommand{\fov}{f_{\rm ov}}
\newcommand{\bracketraw}{\mathcal{B}^{\rm raw}}
\newcommand{\bracketb}{\mathcal{B}}
\newcommand{\bracketsm}{\widetilde{\mathcal{B}}^{\rm raw}}
\newcommand{\bracketstar}{\mathcal{B}^{*}}

\shorttitle{Rotation and Overshoot in ORCHARD}
\shortauthors{Tejada Arevalo}

\begin{document}

\title{Rotation, Entrainment, and Convective Overshoot in the ORCHARD Planetary Evolution Code}

\author{Roberto Tejada Arevalo}

\begin{abstract}
We describe four optional physics modules added to the ORCHARD planetary
evolution code.  Together they provide a flux-based representation of
rotation-modified convection and overshoot mixing.  The first module
(\texttt{rotational\_convection}) implements the rotation-suppressed
mixing-length convective velocity from \citet{Fuentes2023ApJL} and
the Coriolis--Inertial--Archimedean force balance of
\citet{Aurnou2020}.  It multiplies the inviscid \citeauthor{Stevenson1979}
MLT velocity by a smooth interpolation $\Rofunc$ that asymptotes to
$\Roo_{\rm NR}^{1/5}$ in the rapid-rotation limit.  The second
(\texttt{convective\_entrainment}, experimental) translates the
\citet{Linden1975} entrainment hypothesis used in
\citeauthor{Fuentes2023ApJL}'s Eq.~5 into a boundary-localised
effective diffusivity.  However, the resulting one-cell-at-a-time front
advance is dynamically pinned by the Ledoux criterion, and does not
reproduce the dramatic erosion their three-dimensional simulations
predict.  The third module
(\texttt{smooth\_convection\_criterion}) replaces the hard
Schwarzschild/Ledoux gate on the diffusion coefficient with a smooth
sigmoid blend evaluated on a spatially Gaussian-smoothed entropy
bracket.  This addresses a long-standing discretisation artefact that
produces alternating convective/stable cells with diffusion
coefficients spanning eleven orders of magnitude.  The fourth
(\texttt{composition\_overshoot}) introduces an exponentially
decaying composition diffusivity beyond every convective/stable
boundary, in the spirit of \citet{Herwig2000} and \citet{Freytag1996}.
An optional kinetic-energy-flux coupling ties the overshoot
strength to $R^3 = (F_{K,R}/F_{K,\rm NR})$ from the
\citeauthor{Fuentes2023ApJL} energy balance.  We derive each
prescription, document the discretisation and Newton--Raphson
Jacobian additions, and present same-age $5$~Gyr Jupiter
inhomogeneous-evolution comparisons that quantify the effect of each
module on the fuzzy-core gradient.
\end{abstract}

\keywords{planets and satellites: interiors --- planets and satellites:
gaseous planets --- convection --- methods: numerical}

% ======================================================================
\section{Introduction and motivation}\label{sec:intro}
% ======================================================================

The interior structure of Jupiter and Saturn, as constrained by Juno
and ring seismology, is consistent with extended composition gradients
(``fuzzy cores'') stable against convective mixing over the age of the
solar system \citep{Wahl2017,Mankovich2021,Howard2023}.  This is
puzzling.  Mixing-length theory (MLT) gives sub-Gyr
homogenisation timescales for the relevant Rayleigh numbers, and
three-dimensional hydrodynamic simulations of layered convection have
shown that double-diffusive staircases erode rapidly through
overshoot mixing \citep{Wood2013,Garaud2018,Fuentes2022PhRvF}.
\citet{Fuentes2023ApJL} (hereafter F23) argued that rapid planetary
rotation is a missing ingredient: in the limit
$\Roo \equiv \Vconv / (2\Omega \ell) \ll 1$, Coriolis forces reduce
convective velocities by a factor of $\sim 6$ and the kinetic-energy
flux available for entrainment by $\sim 200$, producing a
$\Roo^{-1/2}$ stretch of the boundary-mixing timescale that may bring
the fuzzy-core erosion time within the age of Jupiter.

ORCHARD is a one-dimensional Lagrangian Henyey/Newton--Raphson
code, so its representation of convection is fundamentally a
flux-based MLT closure rather than a hydrodynamic simulation
\citep{TejadaArevalo2026orchard}.  This document describes four
optional physics modules that incorporate F23-style rotation
suppression and overshoot mixing within ORCHARD's flux framework.
None of the modules attempts to simulate fluid parcels.  Each
instead translates a physical effect into a modification of the
diffusion coefficient $D_b$, the convective velocity $\Vmlt$, or
the convective heat flux $\mathcal{F}_{\rm conv}$.  All four
modules default to disabled so that existing models reproduce the
fiducial flux closure exactly.

The four modules are:
\begin{enumerate}
\item \texttt{rotational\_convection} (\S\ref{sec:rot}) --- the
  rotation-modified MLT velocity of F23 / \citet{Aurnou2020}, with a
  smooth interpolation between rapid-rotation and slow-rotation limits.
\item \texttt{convective\_entrainment} (\S\ref{sec:entrainment},
  experimental) --- the \citet{Linden1975} entrainment hypothesis
  used in F23 Eq.~5, translated into a boundary-localised effective
  diffusivity.
\item \texttt{smooth\_convection\_criterion} (\S\ref{sec:smooth}) ---
  a sigmoid blend of $\Dmlt$ and $\Dmicro$ that smooths the discrete
  Schwarzschild/Ledoux gate.
\item \texttt{composition\_overshoot} (\S\ref{sec:overshoot}) --- an
  exponentially decaying composition diffusivity at every convective
  boundary, with an optional $R^3$ rotation coupling that inherits the
  F23 energy-flux suppression.
\end{enumerate}

We document the physics derivation, discretisation, and Newton--Raphson
Jacobian additions for each module, then present numerical comparisons
on a non-rotating Jupiter inhomogeneous--Gaussian model evolved to $5$
Gyr (\S\ref{sec:results}).

% ======================================================================
\section{Notation and baseline closure}\label{sec:notation}
% ======================================================================

We adopt the same notation as the ORCHARD transport methods document.
Mass-shell boundaries are indexed $b = 0, 1, \ldots, N$ with $b = 0$
the surface and $b = N$ the centre, so the radial coordinate $r_b$
descends with $b$.  Cell-centred quantities such as the
specific entropy $S$, helium mass fraction $Y$, and metal mass
fraction $Z$ are stored at integer cell indices $k$ with
$k \in [0, N-1]$.  We denote the entropy bracket
\begin{equation}
\bracketb{}_b
    = -\frac{T H_P}{\Cp} \frac{\dd S}{\dd r}\bigg|_b
      + \mathrm{(composition\;terms)}
    \label{eq:bracket-defn}
\end{equation}
which is positive in convective zones and (post-softhinge) zero in
stable zones.  In the Ledoux closure the composition terms include
contributions of the form $(\partial S/\partial X_i)_{\rho,P}\,
\partial X_i/\partial r$ that stabilise against compositional
gradients.  The pre-softhinge bracket $\bracketraw_b$ retains its
sign and is generally negative in stable regions.

In the standard inviscid MLT closure used by ORCHARD
\citep[Eq.~11 of][]{TejadaArevalo2026orchard},
\begin{equation}
\Vmlt^2
   = \frac{g\,\delta\,(\nabla - \nabla_{\rm ad})\,\ell^2}{8 \Hp}
   = \frac{g\,\Hp^2}{8\Cp}\,|\bracketb|,
   \label{eq:vmlt}
\end{equation}
with $\ell = \Hp$ the mixing length, $\delta = -(\partial \ln \rho/
\partial \ln T)_P$, and the second equality assumes the entropy-gradient
form of the bracket.  The MLT composition diffusivity follows from
\begin{equation}
\Dmlt = \frac{\Vmlt \Hp}{3}.
\label{eq:Dmlt}
\end{equation}
Composition is then evolved by a sum
\begin{equation}
D_b = \begin{cases}
\Dmlt + \Dmicro & \bracketb{}_b > 0,\\
\Dmicro & \bracketb{}_b \le 0,
\end{cases}
\label{eq:Db-default}
\end{equation}
which is the hard step the smooth-convection module replaces.  The
softhinge \citep[\S2.5 of][]{TejadaArevalo2026orchard} clamps
$\bracketraw_b$ to $\ge 0$ before $\Vmlt$ is taken so that $\Vmlt$ is
real-valued; in the Newton solver this introduces an exact
discontinuous transition between $D \sim 10^{10}\,\rm cm^2/s$ and
$D = \Dmicro \sim 0.05\,\rm cm^2/s$ across one boundary.

The angular velocity $\Omega(t)$ is computed in
\texttt{hydrostatic.py:get\_omega} from angular-momentum
conservation \citep[Appendix A of][]{TejadaArevalo2026orchard}.  It is
threaded into \texttt{transport.update\_SYZ} via an
\texttt{omega\_old} keyword argument.  When the master flag
\texttt{[hydrostatic\_equilibrium].rotation = False}, $\Omega = 0$
identically and all rotation-coupled corrections collapse to identity.

% ======================================================================
\section{Rotation-modified convective velocity}\label{sec:rot}
% ======================================================================

\subsection{Physics}

Following F23, in the rapidly rotating limit the convective force
balance is set by the Coriolis, inertial, and Archimedean (CIA) force
ratio rather than by the inertial--Archimedean balance that gives the
\citet{Stevenson1979} non-rotating MLT velocity.  Define the
non-rotating Rossby number per zone boundary,
\begin{equation}
\Roo_{\rm NR} \equiv \frac{\VNR}{2\Omega\Hp},
\label{eq:Ro-NR}
\end{equation}
where $\VNR$ is the standard MLT velocity (Eq.~\ref{eq:vmlt}) using
$\ell = \Hp$.  The CIA balance derived in
\citet[][Eq.~3]{Fuentes2023ApJL} gives in the strongly rotating limit,
\begin{equation}
\frac{\VR}{\VNR}
    \;\sim\; \left(\frac{\VNR}{2\Omega\Hp}\right)^{1/5}
    \;=\; \Roo_{\rm NR}^{1/5}.
\label{eq:VRovervNR-asymp}
\end{equation}
At Jovian deep-envelope conditions ($\VNR \approx 40\,\rm cm/s$,
$\Omega = 2\pi/10\,$hr, $\Hp \approx 10^9\,\rm cm$), $\Roo_{\rm NR}
\approx 10^{-4}$, so $\VR/\VNR \approx 0.16$, reproducing
F23's reported factor-of-six velocity reduction.

In the slow-rotation limit ($\Roo_{\rm NR} \gg 1$) Coriolis forces are
negligible and $\VR \to \VNR$.  We bridge the two limits with a smooth
Stevenson-style interpolation,
\begin{equation}
\Rofunc \equiv \left(\frac{\Roo_{\rm NR}}{1 + \Roo_{\rm NR}}\right)^{1/5},
\label{eq:Rfunc}
\end{equation}
which is algebraically equivalent to $(1 + \Roo_{\rm NR}^{-1})^{-1/5}$
but does not require dividing by a small Rossby number at deep
envelope depths.  The function satisfies
$R \to \Roo_{\rm NR}^{1/5}$ as $\Roo_{\rm NR} \to 0$ and
$R \to 1$ as $\Roo_{\rm NR} \to \infty$.  It therefore matches F23
Eq.~3 in the rapid-rotation regime and the standard MLT in the
slow-rotation regime.

The rotation-modified convective velocity is then
\begin{equation}
\VR = \VNR \cdot \Rofunc,
\label{eq:VR}
\end{equation}
and the convective heat flux and MLT composition diffusivity inherit
the same factor:
\begin{equation}
\mathcal{F}_{\rm conv,R} = R \cdot \mathcal{F}_{\rm conv,NR},
\quad\quad
D_{\rm MLT,R} = R \cdot D_{\rm MLT,NR}.
\end{equation}
For self-consistency, when ORCHARD's Newton--Raphson transport solve
is asked to carry the same internal flux $\mathcal{F}_{\rm conv}$
that the cooling sequence demands, the entropy bracket adjusts as
$|\bracketb|_R / |\bracketb|_{\rm NR} = R^{-2/3}$ (about $3\times$
steepening at Jovian rotation), since the bracket-form flux scales as
$\mathcal{F}_{\rm conv} \propto |\bracketb|^{3/2}$ and the prefactor
$R$ shifts the response.

\subsection{Discretisation and Jacobian}

The implementation lives in
\texttt{transport.rotation\_factors}, called once per Newton iterate
in \texttt{update\_SYZ}.  Given $\VNR$ on every boundary
(reused from the standard MLT computation at line 2670 of
\texttt{transport.py}) and $\Omega$ from
\texttt{hydrostatic\_equilibrium}, we evaluate
\begin{align}
\Roo_{{\rm NR},b} &= \frac{V_{{\rm NR},b}}{2\,\Omega\,H_{P,b}}, \\
R_b   &= \left(\frac{\Roo_{{\rm NR},b}}{1+\Roo_{{\rm NR},b}}\right)^{1/5},
\end{align}
on envelope boundaries $b \in [1, k_{\rm core})$, leaving
mantle/core boundaries with $R_b = 1$ identically.  We then multiply
the convective heat-flux prefactor $\phi_b'$ and the MLT diffusivity
$\Dmlt$ by $R_b$.

The rotation factor depends on the entropy bracket through
$\VNR \propto |\bracketb|^{1/2}$.  The Newton--Raphson Jacobian
therefore acquires a chain-rule contribution.  Using
$\dd\ln R / \dd \ln \Roo = 1 - R^5$ and
$\dd\ln \Roo / \dd \ln |\bracketb| = 1/2$ we have
\begin{equation}
\frac{\dd\ln R}{\dd \ln |\bracketb|} = \tfrac{1}{10}(1 - R^5).
\end{equation}
The flux Jacobian then carries an effective bracket-power factor
\begin{equation}
R_{\rm eff} \equiv p + \tfrac{1}{10}(1 - R^5),
\label{eq:Reff}
\end{equation}
with $p = 3/2$ the unmodified MLT exponent, and the existing
\texttt{d\_flux\_dS} array is rescaled by $R_{\rm eff}/p$.  The MLT
diffusivity Jacobian acquires
\begin{equation}
D_{\rm eff} \equiv \tfrac{1}{2} + \tfrac{1}{10}(1 - R^5),
\label{eq:Deff}
\end{equation}
i.e.\ a multiplicative correction $D_{\rm eff}/(1/2) = 2 - R^5$ on
the existing $\dd \Dmlt/\dd|\bracketb|$ term.  Both factors collapse
to identity when $\Omega = 0$ (so $R_b = 1$ and $R^5 = 1$).

\subsection{Configuration}

Activate via
\texttt{[transport].rotational\_convection = True} and require
\texttt{[hydrostatic\_equilibrium].rotation = True} with a defined
\texttt{period}.  Outside the H/He envelope (i.e.\ $i \ge k_{\rm core}$)
$R_b = 1$ identically.  The viscous-mantle convection regime is
therefore not modified.

% ======================================================================
\section{Boundary entrainment (experimental)}\label{sec:entrainment}
% ======================================================================

\subsection{Physics}

F23's rotation argument relies on a \emph{kinetic-energy-flux} bound
on the rate at which a convective region can entrain composition from
a neighbouring stable layer.  The \citet{Linden1975} entrainment
hypothesis equates the rate of change of potential energy across the
mixing front with the kinetic-energy flux of the convective motions:
\begin{equation}
\tfrac{1}{4}\, \rho_0\, \beta\, g\, |dC_0/dz|\, h^2\, \frac{\dd h}{\dd t}
    \;\approx\; \tfrac{1}{2}\,\rho_0\, \Vconv^3
    \quad\text{(F23 Eq.~5)},
\label{eq:F23-eq5}
\end{equation}
where $h(t)$ is the depth of the convective layer above the gradient,
$\beta = (\partial \ln \rho /\partial C)|_{P,T}$, and the kinetic-energy
flux is $\mathcal{F}_K = \tfrac{1}{2}\rho \Vconv^3$.  Solving for the
boundary-recession velocity,
\begin{equation}
v_{\rm ent} \equiv \frac{\dd h}{\dd t}
    \;\approx\; \frac{2\Vconv^3}{\NCsq h^2},
\quad\quad
\NCsq \equiv \beta g |dC_0/dz|.
\label{eq:vent}
\end{equation}
With $\Vconv = \VR = \VNR R$, $v_{\rm ent}$ inherits an $R^3 \approx
200\times$ rotation suppression, which is the headline F23 result.

\subsection{Discretisation}

Implemented in \texttt{transport.entrainment\_diffusivity}, gated by
\texttt{[transport].convective\_entrainment}.  We approximate
$\NCsq$ on every zone boundary via
\begin{equation}
\NCsq{}_{,b} = g_b \cdot \max\!\left(\beta_Y |\dd Y/\dd r|_b,\; \beta_Z |\dd Z/\dd r|_b\right),
\end{equation}
with order-of-magnitude composition density derivatives
$\beta_Y \approx 0.5$ and $\beta_Z \approx 2$.  The convective layer
thickness $h_b$ is the integral of a smooth conv/stable indicator
$\sigma_b$ from the boundary outward to the surface.  The boundary is
localised by
\begin{equation}
\xi_b = \left|\frac{\dd \sigma_b}{\dd r}\right|,
\end{equation}
and the resulting effective diffusivity is
\begin{equation}
\Dent{}_b
    \;=\; v_{\rm ent,b}\,\xi_b\,\dd r_b\,h_b
    \quad
    \text{(advection $\to$ diffusion equivalence)}.
\label{eq:Dent}
\end{equation}
The factor $\xi_b\,\dd r_b$ is dimensionless ($\sim 1$ at the
boundary, zero elsewhere), and $h_b$ provides the length scale for
front propagation.  Their product gives $\Dent$ units of $\rm cm^2/s$.
A smooth saturation cap $v_{\rm ent} \le \Vconv$ prevents
divergence as $\NCsq \to 0$.  The Jacobian uses Picard lagging on
$v_{\rm ent}$ (computed once per Newton iteration with frozen $S, Y,
Z$), so the cubic chain rule through $\Vconv^3$ is not required.
The implicit $C$ update is handled by the standard diffusion
Jacobian.

\subsection{Why this module is experimental}

In tests on the Jupiter inhomogeneous--Gaussian model with rotation
disabled, $\Dent$ produced a one-cell advance of the convective
boundary (mass coordinate $0.269 \to 0.278$) within the first
$0.5$~Gyr.  The boundary then \emph{stalled}.  The reason is
fundamental to a flux-based code: as soon as $\Dent$ moves a small
amount of heavy material into the cell adjacent to the convective
zone, the Ledoux bracket of that cell becomes negative (the
composition gradient stabilises against thermal convection), the MLT
diffusivity collapses to $\Dmicro$ in the now-Ledoux-stable cell, and
$\xi_b$ relocalises one cell deeper.  The boundary advances cell by
cell, but each step requires a full Ledoux flip of the next cell,
which is a slow process driven by the small $\Dent$ alone---not the
fast, parcel-driven process F23's three-dimensional simulations
report.  By 5~Gyr the same-age comparison shows only a $1$-cell
boundary shift and a $\sim 7\%$ change in $Z$ at the gradient's
upper edge.

The module is preserved in the public release because (i) it
faithfully implements F23 Eq.~5 in a flux framework, and (ii) it
documents the failure mode that motivates the
\texttt{composition\_overshoot} module of \S\ref{sec:overshoot}.
We do \emph{not} recommend enabling
\texttt{convective\_entrainment} in production runs.  The
overshoot module captures the same physics in a way that is not
Ledoux-pinned.

% ======================================================================
\section{Smooth convective criterion}\label{sec:smooth}
% ======================================================================

\subsection{Motivation: the staircase artefact}

Within composition-gradient regions the Newton solver places adjacent
zone boundaries alternately on $\bracketb{}_b > 0$ (convective; full
MLT diffusivity $\sim 10^{10}\,\rm cm^2/s$) and $\bracketb{}_b = 0$
(stable; molecular $\Dmicro \sim 0.05\,\rm cm^2/s$).  The
diffusion coefficient therefore alternates by eleven orders of
magnitude between adjacent cells, producing a discrete-cell
``staircase'' in $D(r)$ and a corresponding staircase in $Z(m)$ at
later times.  The pattern is a discretisation artefact of the
hard step (Eq.~\ref{eq:Db-default}) combined with the softhinge that
clamps $\bracketb \ge 0$ in stable cells.  Smoother brackets do not
remove it, because the Newton solver still finds equilibria with
alternating convective/stable cells.

\subsection{Sigmoid blend on a spatially smoothed bracket}

The \texttt{smooth\_convection\_criterion} module replaces
Eq.~\ref{eq:Db-default} with a smooth blend
\begin{equation}
D_b \;=\; \sigmaconv{}_b\, \widetilde{\Dmlt}_b + \Dmicro,
\label{eq:Db-blend}
\end{equation}
where $\widetilde{\Dmlt}$ uses $|\bracketraw|$ (so it is real-valued
in stable cells where $\bracketraw < 0$) and the smooth indicator is
\begin{equation}
\sigmaconv{}_b
    \;=\; \frac{1}{1 + \exp\!\left(-k\,\bracketsm_b/\bracketstar\right)},
\label{eq:sigma-conv}
\end{equation}
with $k$ the sharpness, $\bracketstar$ a calibration scale, and
$\bracketsm_b$ a \emph{spatially smoothed} version of
$\bracketraw_b$:
\begin{equation}
\bracketsm_b \;=\; \mathcal{G}_{\sigma_{\rm cell}} \star \bracketraw_b,
\label{eq:bracket-smooth}
\end{equation}
where $\mathcal{G}_{\sigma_{\rm cell}}$ is a one-dimensional Gaussian
kernel of width $\sigma_{\rm cell}$ cells (default $2$).  The
spatial smoothing is essential.  Applying the sigmoid alone leaves the
staircase intact, because adjacent cells with brackets
$\{+\delta, 0, +\delta, 0, \ldots\}$ map to
$\{1, \tfrac{1}{2}, 1, \tfrac{1}{2}, \ldots\}$, which is still a
staircase.  Smoothing $\bracketraw_b$ across $\sim 2$ cells averages
the alternating values.  Adjacent cells then acquire nearly equal
$\sigmaconv$ in the gradient region, and the diffusivity transition
becomes graded over $\sim 4$--$8$ cells rather than discontinuous
across one.

If $\bracketstar \le 0$, it is auto-calibrated per timestep as the
$25$th percentile of the positive bracket distribution.  Default
$k = 10$ gives a moderately sharp transition.  Large $k$ recovers the
hard step, small $k$ produces a very gradual blend.

\subsection{Jacobian additions}

The standard $\dd \Dmlt/\dd|\bracketb|$ chain rule used by the existing
Newton solver continues to operate on the
$\sigmaconv\,\widetilde{\Dmlt}$ product through
$\widetilde{\Dmlt} \propto |\widetilde{\bracketraw}|^{1/2}$.  The new
contribution is the sigmoid derivative:
\begin{equation}
\frac{\dd \sigmaconv{}_b}{\dd \bracketraw_b}
    = \frac{k}{\bracketstar}\, \sigmaconv{}_b\,(1 - \sigmaconv{}_b),
\end{equation}
which appears as an additive term in
$\dd D_b/\dd \bracketb$:
\begin{equation}
\frac{\dd D_b}{\dd \bracketb{}_b}
    = \sigmaconv{}_b\,\frac{\dd \widetilde{\Dmlt}_b}{\dd \bracketb{}_b}
      + \widetilde{\Dmlt}_b\,\frac{\dd \sigmaconv{}_b}{\dd \bracketb{}_b}.
\end{equation}
Because the auto-calibrated $\bracketstar$ and the Gaussian kernel
both depend on \emph{all} cells in the gradient region, fully exact
chain-ruling through these two operations would break the banded
block-tridiagonal structure of the Newton Jacobian.  We therefore
treat the kernel and the auto-calibration as Picard-lagged constants
within a single Newton iteration.  Both are recomputed from scratch
once per outer iteration, but are not differentiated.

\subsection{Verification}

In a $5$~Gyr non-rotating Jupiter inhomogeneous--Gaussian comparison,
\texttt{smooth\_convection\_criterion} reduces the bracket
positive$\leftrightarrow$zero flip count in the gradient region
(mass-coordinate window $0.20$--$0.30$) from $3$ to $1$, and the
geometric-mean diffusivity in that window drops from
$\sim 60\,\rm cm^2/s$ (dominated by isolated convective cells with
$D \sim 10^{10}$) to a uniform $\sim 1\,\rm cm^2/s$.  The Z gradient
profile changes from a stair-step (plateaus at $Z = 0.07, 0.13, 0.19,
0.23$) to a smooth slope.  The atmospheric $Z$ at $5$~Gyr changes by
$\sim 7\%$.  The Newton solver requires roughly an order of
magnitude more wall-clock time per timestep with this module enabled.
The wider effective convective zone increases the implicit
diffusion stencil.

% ======================================================================
\section{Composition overshoot}\label{sec:overshoot}
% ======================================================================

\subsection{Physics: Herwig-style overshoot with $F_K$ coupling}

In stellar evolution practice the most successful flux-based
representation of overshoot mixing is an exponentially decaying
diffusivity beyond the convective boundary
\citep{Freytag1996,Herwig2000}:
\begin{equation}
\Dov(r)
    \;=\; \fov\,\Dmlt(r_b)\,
          \exp\!\left(-\frac{|r - r_b|}{f_{H_P}\Hp(r_b)}\right),
\label{eq:Dov}
\end{equation}
where $r_b$ is the boundary radius, $\fov$ is a calibration prefactor,
and $f_{H_P}$ is the penetration scale in pressure-scale-height units.
Unlike $\Dent$ from \S\ref{sec:entrainment}, $\Dov$ is non-zero in
\emph{stable} cells regardless of bracket sign.  The gradient
region therefore erodes continuously from above.  It does not become
Ledoux-pinned after a single cell of advance.

To couple this prescription to the F23 rotation suppression we
multiply $\fov$ at each boundary by $R_b^3$, where $R_b$ is the
\texttt{rotational\_convection} factor.  The motivation is that
$\fov$ is empirically calibrated against three-dimensional simulations
that themselves measure the kinetic-energy flux available for
overshoot.  F23 Eqs.~3--4 give
$\mathcal{F}_{K,R}/\mathcal{F}_{K,\rm NR} \approx \Roo^{1/2} \sim
R^3$ in the rapidly rotating limit.  A single multiplicative $R_b^3$
factor on $\fov$ therefore inherits the F23 suppression directly,
without any double-counting of the rotation effect:
\begin{equation}
\fov^{\rm rot}_b \;=\; \fov \cdot R_b^3,
\end{equation}
collapsing to $\fov$ identically when $\Omega = 0$.  Note that $\Dmlt$
itself already carries $R$ from the
\texttt{rotational\_convection} module of \S\ref{sec:rot}.  The
combined rotation suppression of $\Dov$ is therefore
$R \cdot R^3 = R^4 \sim 5 \times 10^{-4}$ at Jovian conditions.  We
treat this as a feature: the F23 entrainment timescale scales as
$R^{-1/2}$, but since the entrainment-mixing timescale also depends
on the convective velocity that does the mixing
(via $\Dmlt$), the additional $R$ factor is physically expected.

\subsection{Discretisation}

Implemented in \texttt{transport.composition\_overshoot\_diffusivity},
gated by \texttt{[transport].composition\_overshoot}.  The helper
\begin{enumerate}
\item Identifies every convective $\to$ stable transition by finding
  boundaries where $\sigmaconv{}_b > 1/2$ and $\sigmaconv{}_{b+1} < 1/2$
  (going inward).  $\sigmaconv$ is taken from the
  \texttt{smooth\_convection\_criterion} module if active, or freshly
  computed from a tanh of $\bracketb$ if not.
\item For each transition, applies Eq.~\ref{eq:Dov} to all envelope
  boundaries with $b \ge b_{\rm transition}$, using
  $D_{\rm anchor} = \Dmlt(r_b)$ at the transition cell.
\item When \texttt{composition\_overshoot\_rotation\_couple = True},
  multiplies $\fov$ at each transition by $R_b^3$.
\item Combines overlapping transitions with $\max$ rather than $+$ to
  avoid double-counting overshoot from two boundaries that share the
  same stable region.
\end{enumerate}
The resulting $\Dov{}_b$ is added to both $D_b$ and $D_{Z,b}$.
Helium and metals therefore share the same overshoot diffusivity.  The
Jacobian uses Picard lagging on $\Dov{}_b$.  We treat it as a known
coefficient within each Newton iteration, and let the standard
diffusion Jacobian handle the implicit $C$ update.

\subsection{Calibration}

Stellar overshoot calibrations \citep[e.g.][]{Herwig2000,Pedersen2018}
typically give $\fov \approx 0.01$--$0.1$.  Those are anchored to
$\Dmlt \sim 10^{15}\,\rm cm^2/s$ in stellar cores.  ORCHARD's
gas-giant envelope has $\Dmlt \sim 10^{10}\,\rm cm^2/s$, so the same
$\fov = 0.01$ produces $\Dov \sim 10^{8}\,\rm cm^2/s$ at the boundary
and homogenises the gradient region in $\lesssim 1\,\rm Myr$
(catastrophically too fast).  We default to $\fov = 10^{-10}$,
giving $\Dov \sim 1\,\rm cm^2/s$ at the boundary and Gyr-timescale
erosion that matches F23 expectations for an unrotated Jupiter.  The
default $f_{H_P} = 0.5$ matches stellar overshoot calibrations.
Users should tune $\fov$ for non-Jovian targets.

\subsection{Verification}

In a $5$~Gyr non-rotating Jupiter inhomogeneous--Gaussian comparison,
\texttt{composition\_overshoot} with $\fov = 10^{-10}$ produces:
\begin{itemize}
\item A clean exponential $\Dov(r)$ tail visible from the convective
  boundary at $\sim 3\,\rm cm^2/s$, decaying to $\sim 0.6\,\rm cm^2/s$
  about ten cells deeper.
\item Advance of the effective convective edge from mass coordinate
  $0.269$ (baseline) to $\sim 0.18$ (overshoot) by $5$~Gyr; the
  convective region has consumed the upper half of the fuzzy core.
\item Homogenisation of $Z$ to a uniform value $\approx 0.089$ from
  the surface down to mass coordinate $\sim 0.18$, then a continuous
  gradient deeper.  Without overshoot the $Z$ profile is a clear
  staircase from $0.07$ at the envelope plateau through plateaus at
  $0.13, 0.19, 0.23$ before reaching $\sim 0.30$ at the deeper core;
  with overshoot the staircase is gone and the gradient extent is
  reduced.
\end{itemize}
With \texttt{rotational\_convection = True} and
\texttt{composition\_overshoot\_rotation\_couple = True}, the
$R_b^3 \sim 4\times 10^{-3}$ factor at Jovian Rossby number reduces
$\Dov{}_{\rm max}$ from $\sim 3\,\rm cm^2/s$ to $\sim 0.012\,\rm cm^2/s
\ll \Dmicro$, so the overshoot becomes ineffective and the fuzzy core
is preserved on the planet's age.  This is the smoking-gun
\citeauthor{Fuentes2023ApJL} prediction: rotation specifically renders
overshoot-driven mixing inert, even though the same prescription
applies to both rotating and non-rotating planets.

% ======================================================================
\section{Numerical comparison and combined behaviour}\label{sec:results}
% ======================================================================

We evolve a single Jupiter inhomogeneous--Gaussian initial profile
($M = 1\,\Mjup$, $S_{\rm ini} = 8.2\,k_B/\mathrm{baryon}$,
$Y_{\rm ini} = 0.277$, Helled-Stevenson primordial $Z$ profile,
\texttt{cd} H/He EOS, AQUA water EOS, c23 atmosphere, helium rain
enabled, $\alpha_{\rm rain} = 0.1$) for three configurations all with
$\Omega = 0$:
\begin{enumerate}
\item Baseline (default ORCHARD physics, no smooth\_conv, no overshoot,
  no entrainment).
\item Baseline + \texttt{smooth\_convection\_criterion}.
\item Baseline + \texttt{composition\_overshoot} ($\fov = 10^{-10}$).
\end{enumerate}
All three runs use the same grid ($N = 300$), same time-stepping, and
reach an identical final age of $4.9996\,\rm Gyr$.  At the
final age the three $Z$ profiles differ as follows:

\begin{table}[!h]
\centering
\caption{$Z$ profile at $5$~Gyr in the gradient region.  Same physics
  except the indicated module.}
\label{tab:z-profiles-5gyr}
\begin{tabular}{cccc}
\hline
$m/M$ & baseline & smooth\_conv & overshoot \\
\hline
0.30 & 0.069 & 0.064 & 0.089 \\
0.27 & 0.132 & 0.103 & 0.089 \\
0.25 & 0.132 & 0.161 & 0.089 \\
0.22 & 0.186 & 0.234 & 0.089 \\
0.20 & 0.230 & 0.257 & 0.089 \\
0.17 & 0.275 & 0.290 & 0.140 \\
0.15 & 0.304 & 0.312 & 0.224 \\
0.10 & 0.354 & 0.354 & 0.337 \\
0.05 & 0.381 & 0.381 & 0.378 \\
\hline
\end{tabular}
\end{table}
The baseline column shows the canonical staircase
($Z \approx 0.13$ plateau in mass coordinate $0.25$--$0.27$, jumping
to $\approx 0.19$ at $0.22$, etc.).  The smooth\_conv column shows
the same overall gradient with the staircase smoothed into a slope.
The overshoot column shows that the upper half of the gradient
region has been mixed up into the envelope.  There $Z = 0.089$
uniformly from the surface down to $m/M \approx 0.18$, with a
continuous gradient below.  The atmospheric $Z$ enrichment at $5$~Gyr is
$\sim 30\%$ for the overshoot model, modest for smooth\_conv, and
nominal for the baseline.

The same module set with rotation enabled produces qualitatively
different behaviour: the \texttt{rotational\_convection} module alone
steepens the bracket and modestly inflates the planet, but does not
move the fuzzy-core boundary (the Ledoux criterion still pins the
convective edge).  \texttt{composition\_overshoot} with
\texttt{composition\_overshoot\_rotation\_couple = True} drops
$\fov$ effectively by $R^3 \approx 0.004$, suppressing the
overshoot diffusivity below $\Dmicro$ and preserving the fuzzy core
to the planet's age.  The Theory-of-Figures cost of rotating runs
makes a head-to-head $5$~Gyr rotcvn+overshoot test impractical in
a single session.  The analytic argument nevertheless follows directly
from the $R^3$ coupling and the per-boundary $\fov^{\rm rot}_b$ used in
the implementation.

% ======================================================================
\section{Summary}\label{sec:summary}
% ======================================================================

The four modules described above implement a flux-based representation
of rotation-modified convection and overshoot mixing in ORCHARD.
Among the four, \texttt{rotational\_convection},
\texttt{smooth\_convection\_criterion}, and
\texttt{composition\_overshoot} are recommended for production use in
giant-planet evolution; the \texttt{convective\_entrainment} module is
preserved as a faithful translation of the F23 entrainment hypothesis,
but does not capture the dramatic boundary-front propagation that
requires hydrodynamic parcel physics.  Together, the modules let ORCHARD
reproduce the F23 prediction that rotation specifically can preserve
the fuzzy core of an unrotated planet that would otherwise be eroded
on Gyr timescales.

All four modules respect the master flag
\texttt{[hydrostatic\_equilibrium].rotation}, default to disabled,
and leave the existing flux closure exactly intact when off.  Each
module is independently toggleable.  They are designed to be used in
combination, but do not require it.

\bibliographystyle{aasjournal}
\bibliography{convective_rotation}

\end{document}
